\chapter{Quantities, Units, and Scientific Conventions}
\label{chap:units}
\index{quantities}\index{units}\index{siunitx package}\index{uncertainty}

A value without a clear unit can make a result unusable. A value with an
inconsistent decimal mark or ambiguous uncertainty can make a table difficult
to audit. The \pkg{siunitx} package separates a number from its unit, applies a
consistent convention, and aligns numerical columns.

\begin{learningoutcomes}
  \item Format numbers and scientific notation with \pkg{siunitx}.
  \item Typeset quantities, ranges, uncertainties, angles, and percentages.
  \item Align numerical values at the decimal marker.
  \item Distinguish measurement precision from displayed decimal places.
  \item Apply concise conventions for chemistry, taxonomy, constants, vectors,
        and matrices.
\end{learningoutcomes}

\section{Load siunitx and separate roles}

Load the package in the preamble. Use \cmd{num} for a number and \cmd{qty} for
a number with a unit:

\begin{latexcode}{Complete example: numbers and quantities}
\documentclass{article}
\usepackage{siunitx}
\begin{document}
The count was \num{12500}. The distance was
\qty{12.4}{\milli\metre}, and the rate was
\qty{3.2e-4}{\metre\per\second}.
\end{document}
\end{latexcode}

\cmd{num}\required{12500} groups digits according to the package settings.
\cmd{qty} receives a numerical argument and a unit argument. The unit is built
from semantic commands: \cmd{milli}\cmd{metre}; division is expressed by
\cmd{per}. Input \code{3.2e-4} is formatted as scientific notation.

Do not type an italic \code{m} in math mode to imitate metre. A unit is upright
and its spacing from the number follows a convention.

\section{Ranges and lists}

Use dedicated commands rather than inserting a hyphen between two complete
quantities:

\begin{latexcode}{Ranges and lists}
The interval was \qtyrange{4.2}{7.8}{\milli\metre}.

The counts were \numlist{12;15;19}.

Temperatures ranged from
\qty{-5}{\degreeCelsius} to \qty{18}{\degreeCelsius}.
\end{latexcode}

The semicolon separates items in \cmd{numlist} input. The output punctuation
can adapt to the document language and package configuration.

\begin{pausepredict}
Compare typed text \code{5-10 mg} with
\code{\textbackslash qtyrange\{5\}\{10\}\{\textbackslash milli\textbackslash
gram\}}. Which version gives \LaTeX{} explicit knowledge of both the range and
the unit?
\end{pausepredict}

\section{Uncertainty and significant figures}

Parenthetical uncertainty input places the uncertainty in the final digits:

\begin{latexcode}{Quantities with uncertainty}
The estimated length was \qty{12.4(3)}{\milli\metre}.

The mean difference was
\qty[separate-uncertainty=true]{0.42(15)}{\milli\metre}.
\end{latexcode}

The first value represents \(12.4\pm0.3\) millimetres. The option in the
second command requests an explicit plus-minus form. Confirm the uncertainty
notation used by your field and journal; do not infer measurement quality from
the number of digits alone.

Rounding is an analytical and reporting decision. \pkg{siunitx} can round
consistently, but it cannot decide how many significant figures the evidence
supports. Preserve unrounded values in the analysis and apply a documented
display rule.

\section{Percentages, angles, and degrees}

\begin{latexcode}{Percent and angle examples}
Coverage was \qty{95}{\percent}.

The angle was \ang{23.5}.

The temperature was \qty{20.0}{\degreeCelsius}.
\end{latexcode}

\cmd{ang} accepts an angular value; \cmd{degreeCelsius} is a complete unit.
Do not attach a degree sign manually to a Celsius value when a semantic unit is
available.

\section{Decimal alignment in tables}

An \code{S} column parses numbers and aligns them at the decimal marker:

\begin{table}[htbp]
  \centering
  \caption{Simulated paired differences, shown only to demonstrate decimal
  alignment.}
  \label{tab:siunitx-alignment}
  \begin{tabular}{@{}lS[table-format=-2.2]@{}}
    \toprule
    Pair & {Difference (mm)} \\
    \midrule
    P01 & 0.42 \\
    P02 & -1.30 \\
    P03 & 12.05 \\
    \bottomrule
  \end{tabular}
\end{table}

Braces around the column heading protect it from numerical parsing. The
\code{table-format=2.2} reservation allows up to two integer and two decimal
digits; a minus sign is accommodated when present. Inspect the real range
before choosing a format.

\begin{commonerror}
\textbf{Symptom:} a text entry in an \code{S} column produces a parsing error
or poor alignment.

\textbf{Cause:} the column expects numerical input.

\textbf{Correction:} protect a short heading with braces, or give textual rows
a suitable non-numeric column. Do not disguise missing data as an unexplained
number.
\end{commonerror}

\section{Optional discipline-specific conventions}

\subsection{Chemical formulae and reactions}

The \pkg{mhchem} package provides \cmd{ce} for chemical notation:

\begin{latexcode}{Optional chemistry source}
\usepackage[version=4]{mhchem}

Water is \ce{H2O}, and a reaction may be written as
\ce{2H2 + O2 -> 2H2O}.
\end{latexcode}

This fragment shows preamble and body lines, not a complete document. Check the
current \pkg{mhchem} manual for supported reaction arrows, charges, isotopes,
and the required version option.

\subsection{Species names and sequences}

Use italics for a genus-species binomial as the disciplinary convention
requires: \textit{Arabidopsis thaliana}. A DNA sequence such as
\code{ACGTTGCA} is data-like text; a monospaced font can distinguish its fixed
symbols, but long sequences need line breaking, numbering, and often a
specialised package or supplementary file.

\subsection{Physical constants, vectors, and matrices}

State whether a value is exact by definition or measured, give its unit, and
cite the authoritative source when appropriate. Use one declared vector
convention, such as \(\mathbf{v}\), and a matrix convention such as
\(\mathbf{A}\). Do not alternate bold, arrowed, and plain vectors without
meaning.

\begin{scienceinpractice}
An engineering tolerance \qty{25.0(5)}{\milli\metre}, an environmental
concentration \qty{8.2}{\micro\gram\per\litre}, and a physical speed
\qty{3.0e8}{\metre\per\second} share one source grammar while retaining their
disciplinary meanings.
\end{scienceinpractice}

\begin{goodpractice}
Define units in column headings when every value shares the unit. Do not repeat
the unit in every table cell. Explain missing values, uncertainty, rounding,
and transformations in the caption or a table note.
\end{goodpractice}

\chapterproject

\begin{miniproject}
Load \pkg{siunitx} in the running article. Declare that both methods report
millimetres. Format the example range, mean difference, and uncertainty with
\cmd{qty}. Add a three-row simulated decimal-aligned table as a temporary
methods check; Chapter~\ref{chap:tables} will turn it into a complete
publication table.
\end{miniproject}

\chaptersummary

\pkg{siunitx} formats numbers separately from units and supplies commands for
quantities, ranges, lists, uncertainties, percentages, angles, and scientific
notation. Its \code{S} column aligns parsed numbers. Formatting consistency is
not a substitute for justified significant figures. Optional tools such as
\pkg{mhchem} express discipline-specific notation; taxonomy, sequences,
constants, vectors, and matrices also require declared conventions.

\chaptercheck
\begin{chapterchecklist}
  \item Every reported measurement has an unambiguous unit.
  \item Ranges and uncertainties use semantic commands.
  \item Display rounding follows a stated analytical rule.
  \item Numerical table columns align at the decimal marker.
  \item Discipline-specific notation follows a verified convention.
  \item Simulated quantities remain labelled as demonstrations.
\end{chapterchecklist}

\chapterexercisestitle
\begin{chapterexercises}
  \item Typeset \num{1250000} and \qty{2.5e-6}{\metre}.
  \item Create a range from 4 to 9 milligrams.
  \item Express \(7.2\pm0.4\) centimetres using uncertainty input.
  \item Make an \code{S} column for values from \(-12.30\) to \(8.05\).
  \item Explain why decimal places do not by themselves establish precision.
  \item Typeset a Celsius temperature and a percentage.
  \item Identify the role of \cmd{ce} and the package that supplies it.
  \item Challenge: audit a table from your field for units, rounding,
        uncertainty, and missing-value definitions.
\end{chapterexercises}
